\section{The equations on a projective threefold}\label{sec:compact}

Let $B=\Def(X)$ and let $\mathcal X\to B$ be a proper semiuniversal
family. We use the local models of Section~\ref{sec:local} with the
markings fixed in the introduction. A choice of global generator of
$\omega_X$ fixes the scale of the infinitesimal period comparison
at every point.

\subsection{The period map of a global form}
After shrinking $B$, the family is relatively Gorenstein and its
relative dualizing sheaf is a line bundle. Its restriction to $X$
has a generator $\omega$. The vanishing $H^1(X,\omega_X)=0$ permits
a lift of $\omega$ over the full analytic germ. Indeed a finite free
complex for the proper cohomology of the relative dualizing sheaf
computes cohomology after base change. In a minimal such complex the
degree-zero term has rank one and the adjacent terms vanish, by the
cohomology of the central fibre. A generator of the resulting direct
image lifts $\omega$, and its zero locus has closed image disjoint
from the origin. Shrinking once more yields a relative volume form
$\widetilde\omega$. The same calculation for $\OO_{\mathcal X}$
gives $\pi_*\OO_{\mathcal X}=\OO_B$, where $\pi$ is the family map.
This form of the base change argument, including nonreduced $B$,
is also used in \cite[\S2]{Wue}.

For each $p$ choose an analytic classifying map $b_p:B\to B_{Q_p}$
and an identification of the germ of the global family at $p$ with
the pullback of the toric family. Embed $B$ in a smooth germ $D$.
The graph of $b_p$ in $B_{Q_p}\times D$ is closed, so the above
pullback is a closed subgerm of $\mathcal X_{Q_p}\times D$. There
$\widetilde\omega/\Omega_{Q_p}$ is a unit on the pullback. Lift it to
a unit $g_p$ on the ambient family by lifting a representative in the
local ring. Theorem~\ref{thm:local} supplies over $D$ an isomorphism
\[
  B_{Q_p}\times D\longrightarrow G_p\times D
\]
whose first component is the marked period after link duality. Its
restriction to the graph of $b_p$ depends only on the restricted form:
on each Artin quotient it is the de Rham class of
$\widetilde\omega$, and the analytic local ring of $B$ is separated.
Taking all these restrictions gives $P:B\to G=\prod_pG_p$.

For any Artin pullback the smooth locus over the central open set
$U=X_{\rm reg}$ carries the closed form $\widetilde\omega$.
The relative de Rham resolution of the constant coefficient ring
therefore puts its link classes in the image of $H^3(U)\to
\bigoplus_p H^3(L_p)$ tensored with that ring. By
Proposition~\ref{prop:tangent} this image is $K$ after duality.
Since the linear equations of $K$ vanish on every maximal-ideal
quotient of $\OO_{B,0}$, the analytic map $P$ factors through
$K\times_V G$. Its tangent map has image $K$ and kernel
$H^1(X,\Theta_X)$ by the same proposition, now applied to the
central form $\omega$. A different relative generator rescales the
tuple by a single unit on $B$, which preserves both the linear
subspace and the homogeneous local ideals.

\subsection{Lifting with prescribed periods}
\begin{proof}[Proof of Theorem~\ref{thm:global}]
Choose functions $s_1,\ldots,s_h$ on $B$ vanishing at the origin
whose differentials restrict to a basis of the dual of $\ker dP$.
They give a map
\[
 F=(s,P):B\longrightarrow T=\C^h\times (K\times_V G).
\]
Since each $G_p$ has tangent space $V_p$, one has
$T_0T=\C^h\oplus K$ and $dF$ is an isomorphism. In particular $F$
is a closed immersion of germs. Indeed the pullbacks of ambient
coordinates span the cotangent space of $B$. A basis among them is a
coordinate system in a minimal smooth embedding of $B$, by the
analytic inverse function theorem.

Let $A'\to A$ be a small extension with kernel $I$ and let $b\in B(A)$.
Fix $y'\in T(A')$ above $F(b)$. Lift the image of $b$ in the
smooth ambient germ $D$ to $D(A')$. With this auxiliary parameter
fixed, the $p$th component of the prescribed period in $y'$ pulls
back under the inverse local period isomorphism to a point of
$B_{Q_p}(A')$ above $b_p(b)$. This proves local liftability of the
given deformation simultaneously at all of its singular points.
The lifting hypothesis now supplies a global lift. A semiuniversal
family induces a smooth morphism from its formal base to the
deformation functor. Hence this lift is represented by a point
$b'\in B(A')$ reducing to the specified point $b$.

The two homomorphisms $\OO_{T,0}\to A'$ defined by $y'$ and
$F(b')$ have the same reduction to $A$. Their difference is a
derivation at the closed point with values in $I$, because
$\mathfrak m_{A'} I=0$. Let $v\in T_0B\otimes I$ be the inverse
image of this difference under $dF\otimes I$. Adding $v$ to the
homomorphism defining $b'$ again gives a ring homomorphism, and
its image under $F$ is $y'$. This proves formal smoothness of the
closed immersion. Applying the lifting to the successive
maximal-ideal quotients of $\OO_{T,0}$ gives a left inverse to the
completed quotient homomorphism, hence its injectivity.
Faithful flatness of completion gives the asserted analytic
isomorphism.
\end{proof}

\begin{remark}
The lifting argument also describes what follows from local
liftability alone. Suppose $B$ classifies a semiuniversal deformation
with isolated singularities and the lifting hypothesis of
Theorem~\ref{thm:global} holds. Write $D_0$ for the product of its
local hulls, $\lambda:T_0B\to T_0D_0$ for tangent localization, and
$h_0,r$ for the dimensions of its kernel and cokernel. In a minimal
smooth embedding of $D_0$ with ideal $I_0$, the global base admits
a presentation
\[
 \OO_{B,0}\simeq \C\{s_1,\dots,s_{h_0},t_1,\dots,t_{\dim T_0D_0}\}
             /(I_0,f_1,\dots,f_r)
\]
with the $df_i(0)$ linearly independent. To see this, embed $B$ in
$\C^{h_0}\times D_0$ by the local classifying maps and functions
dual to the kernel as above. Choose elements of its ideal whose
linear parts span the conormal space in the smooth ambient germ.
Their zero locus in $\C^{h_0}\times D_0$ has the same tangent space
as $B$ and any point of this locus gives local lifts via its
projection to $D_0$. The same small-extension argument proves
equality of the germs. For the toric points, the periods give the
$f_i$ as linear equations on $V$, independent of the smooth
coordinates.
\end{remark}
